\chapter{Recurrences}\label{ch:rec}
\index{recurrence}

A recurrence describes a sequence by its own past --- the natural output
of the case splits from Chapter~8 and Chapter~3. This chapter is about
turning that description into either a closed form (characteristic
equations, generating functions) or at least a growth order. After this
chapter you can set up a recurrence from a problem's own structure,
solve the linear homogeneous ones, and predict how fast the answer
grows.

% from-zero:begin
\begin{rulebox}[Before you start]
Chapters~3 (induction) and~8--9 (counting, generating functions). \emph{Recap:} a recurrence defines each term from earlier ones, like $F_n=F_{n-1}+F_{n-2}$; a closed form gives $F_n$ directly.
\end{rulebox}
\begin{wins}[After this chapter you can]
\begin{itemize}
\item guess a closed form from the first few terms and prove it by induction;
\item solve linear recurrences with the characteristic equation;
\item set up a recurrence for a tiling or counting problem;
\item estimate how fast a sequence grows.
\end{itemize}
\end{wins}
\begin{principle}[How to think about this]
A recurrence tells you how to get the next term from earlier ones, which is exactly how you would explain the pattern to a friend. The goal is to turn that step-by-step description into a direct formula. Guess from small cases, then \emph{prove} the guess by induction; do not trust a pattern until it is proved.
\end{principle}
% from-zero:end

\section{Recurrence vs closed form}\label{sec:recvs}
\index{closed form}
A \emph{recurrence relation} gives $a_n$ in terms of earlier terms plus
base cases; a \emph{closed form} gives $a_n$ as a function of $n$
directly. The recurrence is usually what the problem gives you for
free; the closed form is what you earn by solving it. Both are honest
answers --- but a closed form that was never verified against the
recurrence is just a conjecture.

\section{Guessing, then proving}\label{sec:guess}
\index{induction!closing a guess}
The standard pipeline: unroll or compute a few values, conjecture a
closed form, then prove it by induction --- the guess supplies the
formula, induction supplies the licence.

\begin{example}\label{ex:hanoi}
Solve the Towers of Hanoi recurrence $H_n = 2H_{n-1} + 1$, $H_1 = 1$.
\par\smallskip
\textbf{Strategy.} Compute four values first --- never conjecture from
one data point. $1, 3, 7, 15$ is $2^n - 1$; the recursion tree
(Figure~\ref{fig:hanoi}) shows \emph{why}: each level doubles, and the
$+1$ is the single real move of the biggest disk.
\par\smallskip
\begin{solution}
Claim $H_n = 2^n - 1$. Base: $H_1 = 1 = 2^1 - 1$. Step:
$H_n = 2H_{n-1} + 1 = 2(2^{n-1} - 1) + 1 = 2^n - 1$ by the induction
hypothesis. Unwrapping gives the same sum directly:
$H_n = 1 + 2 + 4 + \cdots + 2^{n-1} = 2^n - 1$.
\end{solution}
\end{example}

\begin{figure}[ht]\centering
% Figure: Hanoi recursion tree depth 4 with call-count labels.
% Source for Figure~\ref{fig:hanoi}. TikZ only; regenerable by compiling main.tex.
\begin{tikzpicture}[
  nd/.style={circle,draw=rulegray,inner sep=2pt,minimum size=7mm,font=\scriptsize},
  ed/.style={draw=rulegray},
  lab/.style={font=\footnotesize\sffamily,text=slate},
]
\node[nd] (r) at (0,0) {$H_4$};
\node[nd] (a) at (-2.2,-1.1) {$H_3$};
\node[nd] (b) at (2.2,-1.1) {$H_3$};
\draw[ed] (r) -- (a); \draw[ed] (r) -- (b);
\node[nd] (a1) at (-3.2,-2.2) {$H_2$};
\node[nd] (a2) at (-1.2,-2.2) {$H_2$};
\node[nd] (b1) at (1.2,-2.2) {$H_2$};
\node[nd] (b2) at (3.2,-2.2) {$H_2$};
\draw[ed] (a) -- (a1); \draw[ed] (a) -- (a2);
\draw[ed] (b) -- (b1); \draw[ed] (b) -- (b2);
\foreach \p/\q in {a1/a11,a1/a12,a2/a21,a2/a22,b1/b11,b1/b12,b2/b21,b2/b22} {}
\foreach \i/\x in {0/-3.6,1/-2.8,2/-1.6,3/-0.8,4/0.8,5/1.6,6/2.8,7/3.6}
  \node[nd,minimum size=5mm] at (\x,-3.3) {$H_1$};
\draw[ed] (a1) -- (-3.6,-3.3); \draw[ed] (a1) -- (-2.8,-3.3);
\draw[ed] (a2) -- (-1.6,-3.3); \draw[ed] (a2) -- (-0.8,-3.3);
\draw[ed] (b1) -- (0.8,-3.3); \draw[ed] (b1) -- (1.6,-3.3);
\draw[ed] (b2) -- (2.8,-3.3); \draw[ed] (b2) -- (3.6,-3.3);
\node[lab,anchor=east] at (-4.1,0) {depth 0: 1 call};
\node[lab,anchor=east] at (-4.1,-1.1) {depth 1: 2 calls};
\node[lab,anchor=east] at (-4.1,-2.2) {depth 2: 4 calls};
\node[lab,anchor=east] at (-4.1,-3.3) {depth 3: 8 calls};
\node[lab] at (0,-4.2) {$1+2+4+8 = 15$ moves --- the tree doubles per level};
\end{tikzpicture}

\caption{The Hanoi recursion tree, four levels deep ($1+2+4+8$ calls). Notice the shape is the
proof: $1 + 2 + 4 + 8$ calls, one real move per call --- the doubling
you see \emph{is} the $2^n - 1$.}
\label{fig:hanoi}
\end{figure}

\begin{figure}[ht]\centering
% Figure: the guess -> check -> prove loop for recurrences.
% Source for Figure~\ref{fig:guessloop}. TikZ only; regenerable by compiling main.tex.
\begin{tikzpicture}[
  st/.style={draw,minimum height=11mm,text width=30mm,align=center,font=\footnotesize,fill=accent!8},
  lab/.style={font=\footnotesize\sffamily,text=slate,align=center}
]
\node[st] (g) at (0,0) {\textbf{Guess}\\compute $a_1..a_6$,\\spot a pattern};
\node[st] (c) at (4.6,0) {\textbf{Check}\\predict $a_7$ from the\\formula, verify};
\node[st] (p) at (2.3,-2.0) {\textbf{Prove}\\induction: base hit,\\step inherited};
\draw[->,thick,accent] (g) -- (c) node[midway,above,lab] {\scriptsize honest trial};
\draw[->,thick,accent] (c) -- (p) node[midway,right,lab,xshift=1mm] {\scriptsize survived};
\draw[->,thick,rulegray,dashed] (p) -- (g) node[midway,left,lab,xshift=-1mm] {\scriptsize failed? refine};
\node[lab] at (2.3,-3.3) {The loop only exits at \emph{Prove} --- a guess\\
that merely fits the data is data, not a theorem.};
\end{tikzpicture}

\caption{The honest pipeline: a computed guess is \emph{evidence}, the
check is a filter, and induction is what converts pattern into theorem.
Skipping the loop is how false patterns get published.}
\label{fig:guessloop}
\end{figure}

\section{Characteristic equations}\label{sec:chareq}
\index{characteristic equation}
\index{linear recurrence!homogeneous}
For $a_n = p\,a_{n-1} + q\,a_{n-2}$, guess $a_n = r^n$ (geometric
ansatz --- geometric sequences are the only ones where each term is a
constant ratio of the previous two, and the recurrence is asking
exactly that question). Substituting gives $r^n = p\,r^{n-1} + q\,r^{n-2}$,
i.e.\ $r^2 = pr + q$ --- the \emph{characteristic equation}.

\begin{example}\label{ex:chareq}
Solve $a_n = 5a_{n-1} - 6a_{n-2}$ with $a_0 = 2$, $a_1 = 5$.
\par\smallskip
\textbf{Strategy.} The ansatz $r^n$ reduces the recurrence to algebra;
two roots mean the general solution is a combination, and the base
cases pin down the weights.
\par\smallskip
\begin{solution}
$r^2 - 5r + 6 = 0$ gives $r = 2, 3$ (Figure~\ref{fig:chareq}).
$a_n = A\cdot 2^n + B \cdot 3^n$; $a_0 = A + B = 2$, $a_1 = 2A + 3B = 5$
gives $A = B = 1$, so $a_n = 2^n + 3^n$. Check: $a_2 = 5\cdot 5 - 6\cdot
2 = 13 = 4 + 9$.
\end{solution}
\end{example}

\begin{figure}[ht]\centering
% Figure: characteristic-equation decision diagram for order-2 recurrences.
% Source for Figure~\ref{fig:chareq}. TikZ only; regenerable by compiling main.tex.
\begin{tikzpicture}[
  bx/.style={draw=rulegray,rounded corners=3pt,inner sep=4pt,align=center,font=\footnotesize},
  dm/.style={draw=accent,inner sep=3pt,align=center,font=\footnotesize},
  ar/.style={->,thick,draw=ink},
  lab/.style={font=\footnotesize\sffamily,text=slate},
]
\node[bx,text width=3.1cm] (r) at (0,0) {$a_n = p\,a_{n-1} + q\,a_{n-2}$:\\try $a_n = r^n$};
\node[bx,text width=3.1cm] (e) at (0,-1.5) {$r^2 = p\,r + q$\\(characteristic equation)};
\node[dm] (c) at (0,-3.0) {roots?};
\node[bx,text width=3.4cm] (d) at (-2.9,-4.5) {distinct $r_1 \ne r_2$:\\$a_n = A r_1^n + B r_2^n$};
\node[bx,text width=3.4cm] (dd) at (2.9,-4.5) {double root $r$:\\$a_n = (A + Bn)\,r^n$};
\draw[ar] (r) -- (e);
\draw[ar] (e) -- (c);
\draw[ar] (c) -- (d) node[midway,above,sloped,lab] {two};
\draw[ar] (c) -- (dd) node[midway,above,sloped,lab] {one};
\node[lab] at (0,-5.6) {$A, B$ are fixed by the two base cases};
\end{tikzpicture}

\caption{The characteristic-equation flowchart. Notice a double root
doesn't halve the answer --- it buys an extra $n$ factor:
$(A + Bn)r^n$, because $n r^n$ is a second independent solution.}
\label{fig:chareq}
\end{figure}

\section{Fibonacci, tilings, and derangements}\label{sec:fib}
\index{Fibonacci}\index{tiling!by recurrence}

\begin{example}\label{ex:tiling}
Count domino tilings of a $2 \times n$ strip.
\par\smallskip
\textbf{Strategy.} Split on the \emph{last} tile's shape, not the first
--- the suffix of a tiling is a shorter tiling, which is exactly the
form a recurrence needs.
\par\smallskip
\begin{solution}
Figure~\ref{fig:tiling}: tilings of $2 \times n$ end either in a
vertical domino (the rest is a $2 \times (n-1)$ tiling, $a_{n-1}$ ways)
or in a horizontal pair ($a_{n-2}$ ways). So $a_n = a_{n-1} + a_{n-2}$
with $a_1 = 1$, $a_2 = 2$ --- the Fibonacci numbers shifted, and
Figure~\ref{fig:ratio} shows the ratio converging to
$\varphi = \frac{1+\sqrt5}{2}$. Binet's formula
$F_n = \frac{\varphi^n - \psi^n}{\sqrt5}$ (with
$\psi = \frac{1-\sqrt5}{2}$) is exactly the two-root case of
\S\ref{sec:chareq}.
\end{solution}
\end{example}

\begin{figure}[ht]\centering
% Figure: 2 x n domino tiling split on the last tile's shape.
% Source for Figure~\ref{fig:tiling}. TikZ only; regenerable by compiling main.tex.
\begin{tikzpicture}[
  gd/.style={draw=rulegray},
  tile/.style={draw=accent,very thick},
  lab/.style={font=\footnotesize\sffamily,text=slate},
]
% case 1: vertical domino at end  -> a_{n-1}
\draw[gd] (0,0) rectangle (4,1);
\draw[gd] (0,0.5) -- (4,0.5);
\draw[tile] (3.5,0) rectangle (4,1);
\node[lab] at (2,-0.4) {ending in a vertical domino: $a_{n-1}$};
% case 2: two horizontal at end -> a_{n-2}
\draw[gd] (0,-1.4) rectangle (4,-0.4);
\draw[gd] (0,-0.9) -- (4,-0.9);
\draw[tile] (3.0,-0.4) rectangle (4,-0.9);
\draw[tile] (3.0,-0.9) rectangle (4,-1.4);
\node[lab] at (2,-1.9) {ending in a horizontal pair: $a_{n-2}$};
\node[lab] at (2,-2.5) {$a_n = a_{n-1} + a_{n-2}$ --- the split is on the \emph{last} tile};
\end{tikzpicture}

\caption{The tiling split on the last domino. Notice only two endings
are possible and they account for disjoint, exhaustive cases ---
that is what licences $a_n = a_{n-1} + a_{n-2}$.}
\label{fig:tiling}
\end{figure}

\begin{figure}[ht]\centering
% Figure: ratio F_{n+1}/F_n converging to the golden ratio.
% Source for Figure~\ref{fig:ratio}. TikZ only; regenerable by compiling main.tex.
\begin{tikzpicture}[
  lab/.style={font=\footnotesize\sffamily,text=slate},
  ax/.style={draw=rulegray},
]
\draw[ax,->] (0,0) -- (6.2,0) node[below left,font=\footnotesize] {$n$};
\draw[ax,->] (0,0) -- (0,3.4) node[above left,font=\footnotesize] {$F_{n+1}/F_n$};
% phi line at y=2.7 (scale: phi-1.6 ~ )
\draw[accent,dashed] (0,2.7) -- (6.0,2.7) node[right,font=\footnotesize] {$\varphi \approx 1.618$};
% ratios: 1,2,1.5,1.667,1.6,1.625,1.615,1.619 -> y = (r-1.55)*20
\foreach \n/\r in {1/1.0,2/2.0,3/1.5,4/1.667,5/1.6,6/1.625,7/1.615,8/1.619}
{
  \pgfmathsetmacro{\y}{(\r-1.4)*1.9}
  \fill[ink] (\n*0.62,\y) circle (1.6pt);
}
\node[lab] at (3.4,-0.5) {the ratio oscillates and converges --- growth factor of every $a_n=a_{n-1}+a_{n-2}$};
\end{tikzpicture}

\caption{$F_{n+1}/F_n \to \varphi$. Notice the oscillation about the
limit --- the shrinking second root $\psi \approx -0.618$ flips sign
each step, which is the oscillation made visible.}
\label{fig:ratio}
\end{figure}

\begin{example}\label{ex:derrec}
Prove the derangement recurrence $D_n = (n-1)(D_{n-1} + D_{n-2})$.
\par\smallskip
\textbf{Strategy.} Classify derangements by where element $1$ lands:
it goes to one of $n-1$ positions $k$; then ask where element $k$ goes.
\par\smallskip
\begin{solution}
Fix where $1$ lands: $n-1$ choices of position $k$. Now element $k$
either lands in position $1$ (the rest is a derangement of the
remaining $n-2$ elements, $D_{n-2}$ ways) or anywhere else (element $k$
still can't go home --- treat position $1$ as its forbidden spot; the
rest is a derangement of $n-1$ elements, $D_{n-1}$ ways). Total
$(n-1)(D_{n-1} + D_{n-2})$.
\end{solution}
\end{example}

\begin{checkpoint}[Predict before you run]
Listing~\ref{lst:fib} runs the recurrence against Binet's formula in
floating point. Predict: does the closed form stay exactly right in
floats, and if not, where does it first break?
\end{checkpoint}

\begin{lstlisting}[language=Python,caption={\texttt{fib\_closed.py}:
recurrence vs Binet --- watch the floats (script in
\texttt{code/fib\_closed.py}).},label={lst:fib}]
def fib_binet(n):
    return (PHI ** n - PSI ** n) / sqrt(5)
\end{lstlisting}
Real output:
\begin{lstlisting}[numbers=none]
n | recurrence | Binet (float) | round | ok?
 10 | 55 | 5.500000e+01 | 55 | True
 20 | 6765 | 6.765000e+03 | 6765 | True
 40 | 102334155 | 1.023342e+08 | 102334155 | True
 50 | 12586269025 | 1.258627e+10 | 12586269025 | True
 60 | 1548008755920 | 1.548009e+12 | 1548008755920 | True
 70 | 190392490709135 | 1.903925e+14 | 190392490709135 | True
 71 | 308061521170129 | 3.080615e+14 | 308061521170130 | False
 72 | 498454011879264 | 4.984540e+14 | 498454011879265 | False
first float failure at n = 71 (double precision; exact answer needs ints)
\end{lstlisting}

Exact in real arithmetic, not in floats: Binet drifts by $n = 71$ in
double precision. The closed form is a theorem about reals; the
recurrence on integers never lies.

\section{Solving by generating functions}\label{sec:recgf}
\index{generating function!solving recurrences}
Chapter~9's conversion, applied: for Fibonacci let
$F(x) = \sum F_n x^n$. The recurrence $F_n = F_{n-1} + F_{n-2}$
multiplied out gives $F(x) = x + xF(x) + x^2F(x)$, so
$F(x) = \frac{x}{1 - x - x^2}$; partial fractions on the denominator's
roots $\varphi, \psi$ recover Binet. The GF turns ``solve a
recurrence'' into ``solve one equation, then do algebra''.

\begin{example}\label{ex:gf3n}
Solve $a_n = 3a_{n-1} + 2$, $a_0 = 1$, by generating functions.
\par\smallskip
\textbf{Strategy.} One conversion: write $A(x) = \sum a_n x^n$, multiply
the recurrence by $x^n$ and sum $n \ge 1$, then solve the linear
equation for $A$ and extract coefficients from a partial-fraction split.
\par\smallskip
\begin{solution}
$\sum_{n\ge1} a_n x^n = 3\sum_{n\ge1} a_{n-1} x^n + 2\sum_{n\ge1} x^n$
gives $A - 1 = 3xA + \frac{2x}{1-x}$, so
$A(x) = \frac{1+x}{\,(1-x)(1-3x)} = \frac{2}{1-3x} - \frac{1}{1-x}$.
Reading coefficients: $a_n = 2\cdot 3^n - 1$. Check: $a_1 = 5 = 2\cdot3-1$
and $a_2 = 17 = 2\cdot9-1$; the recurrence holds since
$3(2\cdot3^{n-1}-1)+2 = 2\cdot 3^n - 1$.
\end{solution}
\end{example}

\section{A feel for growth orders}\label{sec:growth}
\index{growth order}
You can often read the growth order off the recursion tree without
solving: Hanoi splits one call into two calls each one size smaller ---
$2^n$. A recurrence like $T(n) = 2T(n/2) + n$ (merge sort) splits into
two half-size problems; the tree has $\log_2 n$ levels of $n$ work
each, so $T(n) \approx n \log_2 n$. Chapter-17's algorithms live in
this vocabulary; the discipline is to see the \emph{tree} before
reaching for the \emph{formula}.

\begin{figure}[ht]\centering
% Figure: recursion tree of T(n) = 2T(n/2) + n — log n levels, n work each.
% Source for Figure~\ref{fig:rectree}. TikZ only; regenerable by compiling main.tex.
\begin{tikzpicture}[
  nd/.style={circle,draw,minimum size=6.6mm,font=\tiny,fill=accent!8},
  lab/.style={font=\footnotesize\sffamily,text=slate,align=center},
  rw/.style={font=\footnotesize\sffamily,text=accent,align=right}
]
\node[nd] (r) at (0,0) {$n$};
\node[nd] (a) at (-1.6,-0.9) {$n/2$};
\node[nd] (b) at (1.6,-0.9) {$n/2$};
\node[nd] (a1) at (-2.7,-1.9) {$n/4$};
\node[nd] (a2) at (-0.5,-1.9) {$n/4$};
\node[nd] (b1) at (0.5,-1.9) {$n/4$};
\node[nd] (b2) at (2.7,-1.9) {$n/4$};
\draw[->] (r) -- (a); \draw[->] (r) -- (b);
\draw[->] (a) -- (a1); \draw[->] (a) -- (a2);
\draw[->] (b) -- (b1); \draw[->] (b) -- (b2);
\node[lab] at (-0.2,-2.6) {$\vdots$\quad $\log_2 n$ levels \quad $\vdots$};
\node[nd] (l1) at (-1.4,-3.4) {$1$};
\node[nd] (l2) at (0,-3.4) {$1$};
\node[nd] (l3) at (1.4,-3.4) {$1$};
\node[lab] at (0,-4.0) {\dots\quad$n$ leaves\quad\dots};
% per-level work labels on the right
\node[rw] at (4.5,0) {$+n$};
\node[rw] at (4.5,-0.9) {$+n$};
\node[rw] at (4.5,-1.9) {$+n$};
\node[rw] at (4.5,-3.4) {$+n$};
\node[lab] at (4.5,-4.35) {$T(n) = n\cdot\log_2 n + n$};
\end{tikzpicture}

\caption{The recursion tree of $T(n) = 2T(n/2) + n$: $\log_2 n$ levels,
$n$ units of work per level. The closed form is read off the shape ---
$n\log_2 n$, before any formula is invoked.}
\label{fig:rectree}
\end{figure}

\begin{exercise}[easy]\label{exr:rec1}
Unroll $b_n = 3b_{n-1}$, $b_0 = 2$; closed form? Verify $b_4$.
\end{exercise}
\begin{exercise}[medium]\label{exr:rec2}
Solve $a_n = 4a_{n-1} - 4a_{n-2}$, $a_0 = 1$, $a_1 = 4$ (double root ---
use Figure~\ref{fig:chareq}'s second branch).
\end{exercise}
\begin{exercise}[medium]\label{exr:rec3}
Set up the recurrence for climbing an $n$-stair taking 1 or 2 steps,
solve it, and say which sequence this is.
\end{exercise}
\begin{exercise}[medium]\label{exr:rec4}
From \texttt{fib\_closed.py}'s table: why does the float failure show
up only near $n \approx 70$, not earlier? (How many significant digits
does a double carry?)
\end{exercise}

\section{Chapter summary}
\begin{table}[ht]\centering
\begin{tabularx}{\linewidth}{@{}lYY@{}}
\toprule
\tablehead\thead{Shape} & \thead{Method} & \thead{Growth}\\
\midrule
$a_n = c\,a_{n-1} + d$ & unroll / geometric & $c^n$\\
$a_n = pa_{n-1} + qa_{n-2}$ & characteristic eq. & $r_{\max}^n$\\
first-order, side term, $c=1$ & unroll + sum & like the side term\\
split on last object & set up + solve & per the roots\\
linear, any order & generating function & partial fractions\\
halving splits & recursion tree & $n\log n$-ish\\
\bottomrule
\end{tabularx}
\caption{One page: name the shape, pick the method, predict the growth.}
\label{tab:rec}
\end{table}

\section{Common mistakes}
\index{common mistakes}
\begin{itemize}
\item \textbf{Off-by-one base cases.} $a_1 = 1, a_2 = 2$ for tilings is
not $F_1 = 1, F_2 = 1$: the answer is Fibonacci \emph{shifted}. Solve
with the problem's own bases, then reconcile the indices.
\item \textbf{Circular proof.} ``Assume the closed form, plug it into
itself'' proves nothing --- the induction must verify the form
\emph{against the recurrence}, not against itself.
\item \textbf{Formula where the count was the answer.} A closed form
that doesn't evaluate to integers (Binet in floats past $n \approx 70$)
is a liability; when the question wants a count, keep the integers and
treat the formula as a description of growth.
\end{itemize}

\section{Additional exercises}
\begin{exercise}[easy]
Compute $a_5$ for $a_n = 2a_{n-1} + 3$, $a_0 = 1$; then find the closed
form (hint: $a_n + c$ geometric for the right $c$).
\end{exercise}
\begin{exercise}[easy]
Write the recurrence for ``binary strings with no two adjacent 1s'' of
length $n$ and identify the sequence.
\end{exercise}
\begin{exercise}[easy]
Solve $a_n = 7a_{n-1}$, $a_0 = 5$; what is the growth factor?
\end{exercise}
\begin{exercise}[easy]
Solve $a_n = a_{n-1} + n$, $a_0 = 0$ by unwrapping; name the closed
form.
\end{exercise}
\begin{exercise}[medium]
Solve $a_n = 6a_{n-1} - 9a_{n-2}$, $a_0 = 0$, $a_1 = 3$.
\end{exercise}
\begin{exercise}[medium]
Show the $2 \times n$ tiling count with dominoes \emph{and} $2\times2$
squares allowed satisfies $b_n = b_{n-1} + 2b_{n-2}$, and solve it.
\end{exercise}
\begin{exercise}[medium]
Use the derangement recurrence to compute $D_5, D_6$ from
$D_3 = 2, D_4 = 9$, and check $D_6/6!$ against $1/e$.
\end{exercise}
\begin{exercise}[medium]
Find the GF for $a_n = 2a_{n-1} + 3^n$, $a_0 = 1$: set $A(x)$, solve
for it, and extract the closed form by partial fractions.
\end{exercise}
\begin{exercise}[hard]
Prove Binet's formula by induction \emph{and} derive it from
$F(x) = x/(1-x-x^2)$ --- show the two answers coincide term by term.
\end{exercise}
\begin{exercise}[hard]
Extend \texttt{fib\_closed.py} to use \texttt{Decimal} at 50 digits:
does Binet's failure move to a larger $n$, and does \emph{exact}
arithmetic (algebraic form) ever fail? Explain.
\end{exercise}
\begin{exercise}[easy]
Solve $a_n = 2a_{n-1} + 1$, $a_0 = 0$, three ways: unrolling, the
guess-check-prove loop, and the GF of Example~\ref{ex:gf3n}.
\end{exercise}
\begin{exercise}[medium]
Solve $a_n = a_{n-1} + n$, $a_0 = 1$ (arithmetic forcing term ---
which catalog GF produces $n$?).
\end{exercise}
\begin{exercise}[medium]
A path on a staircase of $n$ stones uses steps of $1$ or $3$. Set up
and solve the recurrence for $n \le 9$; where does the characteristic
equation's missing root enter?
\end{exercise}
\begin{exercise}[medium]
Draw the recursion tree of $T(n) = 3T(n/2) + n$ and read off the
growth order without a formula. Compare against $2T(n/2)+n$ ---
what does the branching factor do to the level sums?
\end{exercise}
\begin{exercise}[hard]
Solve the nonhomogeneous $a_n = 4a_{n-1} - 4a_{n-2} + 2^n$,
$a_0 = a_1 = 1$ (resonance: the forcing term is itself a solution of
the homogeneous part --- what replaces the usual ansatz?).
\end{exercise}

\section*{Further reading}
\begin{itemize}
\item R.~Hammack, \emph{Book of Proof}, 3rd ed., Ch.~21 ---
\url{https://richardhammack.github.io/BookOfProof/}
--- free PDF;
induction on sequences and recurrences.
\item E.~Lehman, F.~T.~Leighton, A.~R.~Meyer, \emph{Mathematics for
Computer Science}, Ch.~22 ---
\url{https://courses.csail.mit.edu/6.042/spring18/mcs.pdf}
--- free PDF;
recurrences, characteristic equations and the recursion-tree method.
\item E.~Chen, \emph{Summation} --- generating-function solutions of
linear recurrences:

\url{https://web.evanchen.cc/handouts/Summation/Summation.pdf}
\item Y.~Zhao, \emph{Summations and Combinatorial Identities} ---
olympiad recurrences and the identities they generate:

\url{https://yufeizhao.com/olympiad/comb1.pdf}
\item MIT OpenCourseWare, \emph{6.042J} (Fall 2010) ---
\url{https://ocw.mit.edu/6-042JF10}
--- the recurrences lectures this
chapter follows.
\item Cut-the-Knot, Fibonacci ---
\url{https://www.cut-the-knot.org/}
--- alternate derivations of Binet's
formula for comparison with \S\ref{sec:recgf}.
\end{itemize}

\begin{rulebox}[Watch instead]
\begin{itemize}
\item \textbf{TrevTutor --- ``RECURRENCE RELATIONS''}, 15:24,
\url{https://www.youtube.com/watch?v=eAaP4XaB8hM}
--- Fibonacci and
geometric recurrences, setting up the \S\ref{sec:guess}--%
\S\ref{sec:chareq} pipeline. (Verified 2026-10-11.)
\item \textbf{TrevTutor --- ``HOMOGENEOUS RECURRENCE RELATIONS''}, 25:15,
\url{https://www.youtube.com/watch?v=7mhvA5L7KqY}
--- characteristic
equations worked end to end: \S\ref{sec:chareq}--\S\ref{sec:fib}.
(Verified 2026-10-11.)
\end{itemize}
\end{rulebox}

\chapternotes
The recursion-tree viewpoint and the characteristic-equation pipeline
follow 6.042J's recurrences chapter~\cite{mcs}; the GF-solution section
uses Chen's Summation handout~\cite{chen-summ} as the worked source;
Velleman's recursion section supplied the guess-then-induct
pipeline~\cite{velleman}, and Engel the drill level~\cite{engel}.
Contest items powered by this chapter --- linear recurrences
(AMC~12~2024~A~P21, AMC~12~2024~B~P18, AIME~2021~I~P10), tiling
recurrences (AIME~2025~II~P13, AMC~12~2023~A~P20, AMC~12~2023~A~P25)
and combined counting-recurrence items (AIME~2015~I~P9,
AIME~2023~I~P10) --- are on the AoPS archive~\cite{aops-amc}.
